\section{Cost Model}
\label{sec:cost_model}

This section states our cost model directly, then motivates it with respect
to prior work.

\subsection{Empirical Training Acquisition Cost}
\label{subsec:empirical-cost}

A sampling+learning experiment is parametrized by a configurable system, an
input/workload, a non-functional property $p$, a sampling strategy $s$, a
training size $n$, and a run identifier $r$ (used for stochastic
strategies). Given these, the strategy returns a training sample
$S_{r,s,n}\subseteq C$ with $|S_{r,s,n}|=n$.

The \emph{empirical training acquisition cost} of run $r$ is the sum of
per-configuration acquisition costs over the configurations the strategy
selects:
\begin{equation}
\TrainCost(r,s,n)\;=\;\sum_{c\in S_{r,s,n}} \mathit{cost}(c).
\label{eq:traincost}
\end{equation}
$\mathit{cost}(c)$ is a non-negative scalar that quantifies the resource
consumed when measuring configuration $c$. Returning to the running
example with $\mathit{cost}=$ encoding time: at $n=32$ on x264, $\TrainCost$
is the sum of the $32$ encoding times the strategy decided to pay for, and
Strategy~A and Strategy~B can have very different $\TrainCost$ even though
both have $\SizeCost=32$.

\paragraph*{What can $\mathit{cost}(c)$ be?}
The most common operational instance is wall-clock measurement time --
encoding time for x264, compilation time for LLVM, solving time for DUNE,
runtime for HIPAcc -- and that is the instance we use throughout this
study, because it is the resource that practitioners most often report as
constrained. The model itself is more general, however. Any quantitative
non-negative resource that is consumed by measuring $c$ qualifies as
$\mathit{cost}(c)$:
\begin{itemize}
    \item wall-clock measurement time (the canonical case);
    \item storage, when the operational constraint is accumulated output
    \eg the cumulative size of encoded videos in a transcoding service or
    of compiled binaries in a build farm;
    \item energy or power-integrated consumption, on energy-constrained
    hardware;
    \item bandwidth or I/O volume, on data-centre runs where measurement
    traffic is the bottleneck.
\end{itemize}
The mathematical content of $\TrainCost$ does not change with this choice;
only its operational interpretation does. A reader whose constraint is
storage should read every ``time'' below as ``size'', and the calibration
question and the Pareto question still apply unchanged.

\paragraph*{Target $\neq$ cost.}
The performance property being predicted and the cost dimension are
independent choices. For x264 we predict either encoding time or output
size, and we use encoding time as $\mathit{cost}(c)$ throughout. Predicting
output size while paying time-based cost is one such asymmetric case; an
operator with a storage budget could equally well predict encoding time
while paying size-based cost, and the calibration of $\SizeCost$ against
that $\TrainCost$ would be a different empirical question. We test the
canonical (time-based) instance and discuss generalization in
\Cref{sec:discussion}.

\subsection{Count-Based Baseline}
\label{subsec:count-baseline}

The baseline we compare against is the count of measured configurations:
\begin{equation}
\SizeCost(s,n)\;=\;n.
\label{eq:sizecost}
\end{equation}
This is the cost most often used in the literature
(\Cref{subsec:cost-models}). It is deterministic in $n$, and it is exactly
the abstraction we want to test against $\TrainCost$.

\subsection{Normalized Costs and Calibration Ratio}
\label{subsec:normalized-costs}

Absolute costs vary by orders of magnitude across systems. To compare
across systems we normalize by the cost of measuring the entire valid
configuration space. Define
\begin{equation}
\FullCost\;=\;\sum_{c\in C} \mathit{cost}(c),
\label{eq:fullcost}
\end{equation}
\begin{equation}
\mathit{TrainCostPct}(r,s,n)\;=\;\frac{\TrainCost(r,s,n)}{\FullCost},\quad
\mathit{SizePct}(s,n)\;=\;\frac{n}{|C|}.
\label{eq:pct}
\end{equation}

The size-as-cost abstraction is exact when
$\mathit{TrainCostPct}=\mathit{SizePct}$. We measure how far we are from
that exactness with the calibration ratio
\begin{equation}
\Ratio(r,s,n)\;=\;\frac{\mathit{TrainCostPct}(r,s,n)}{\mathit{SizePct}(s,n)}.
\label{eq:ratio}
\end{equation}
A value of $1$ means the strategy spends a share of the full cost budget
that matches its share of configurations: spending $20\%$ of $|C|$ buys
$20\%$ of $\FullCost$. Values $>1$ mean the strategy is systematically
biased toward expensive configurations; values $<1$ toward cheap ones. We
summarize a strategy by
\begin{equation}
\Bias(s)\;=\;\mathrm{median}_{r,n}\,\Ratio(r,s,n) - 1.
\label{eq:bias}
\end{equation}

\subsection{Composite Objectives}
\label{subsec:composite}

To compare strategies on a single scalar that combines accuracy and cost, we
use a normalized linear combination, instantiated under each cost model:
\begin{align}
J^{\text{size}}_{\lambda}(s,n) &= \widetilde{\MRE}(s,n) + \lambda\cdot\widetilde{\SizeCost}(s,n), \label{eq:Jsize}\\
J^{\text{time}}_{\lambda}(s,n) &= \widetilde{\MRE}(s,n) + \lambda\cdot\widetilde{\TrainCost}(s,n). \label{eq:Jtime}
\end{align}
The tilde denotes min--max normalization within a system/input/property to
prevent scale effects. Rather than committing to a single $\lambda$, we
sweep $\lambda\in\{0.1,0.25,0.5,1,2,4,10\}$, so that the comparison reflects
the full range from accuracy-dominated to cost-dominated trade-offs.

\subsection{Pareto Formulation}
\label{subsec:pareto}

Composite objectives require choosing $\lambda$. To complement the scalar
view with a parameter-free one we also study the two-dimensional MRE--cost
plane. A pair $(s,n)$ is Pareto-efficient under empirical cost if there is
no other pair $(s',n')$ such that
\begin{align*}
\MRE(s',n') &\le \MRE(s,n) \quad\text{and}\\
\TrainCost(s',n') &\le \TrainCost(s,n)
\end{align*}
with at least one strict inequality. The same definition under count-based
cost substitutes $\SizeCost$ for $\TrainCost$. We compare both frontiers in
terms of membership overlap (which $(s,n)$ pairs gain or lose Pareto status
when the cost model changes), the number of Pareto-efficient pairs, and the
front-exclusive pairs under each cost model.

\subsection{Scope}
\label{subsec:scope}

The model above embeds three scope decisions worth stating explicitly,
because they shape what the rest of the paper measures.

\paragraph*{The cost term covers training only.}
A sampling strategy is responsible for the configurations it selects to
train on. The validation set $C_{\text{test}}$ is part of the experimental
protocol; it is the same across strategies and is used only to compute MRE.
We do not include it in $\TrainCost$, so the cost term reflects what the
strategy controls.

\paragraph*{Cost is measured, not assumed.}
$\TrainCost$ is the actual sum of per-configuration costs of the
configurations a strategy selects. We do not assume a per-configuration
unit cost or any closed-form relationship between $n$ and acquisition
cost; we read it off the data.

\paragraph*{Model-building time is reported separately.}
For the learners we use (stepwise linear regression, CART, random forest),
fitting time is orders of magnitude smaller than measurement time when
$\mathit{cost}(c)$ is wall-clock time. We
report it as a separate diagnostic rather than folding it into
$\TrainCost$, so that cost differences across strategies reflect
acquisition effort. \Cref{sec:threats} discusses when this simplification
might be inappropriate.
